% =====================================================================
%  Borrador para arXiv -- preprint
%  Separacion angular maxima en Transformers de vision
%  Compatible con pdfLaTeX, XeLaTeX y LuaLaTeX (Overleaf)
%  estado: borrador a falta de resultados experimentales
% =====================================================================
\documentclass[11pt,a4paper]{article}

% --- motor y codificacion -------------------------------------------
\usepackage{iftex}
\ifPDFTeX
  \usepackage[utf8]{inputenc}
  \usepackage[T1]{fontenc}
\else
  \usepackage{fontspec}
\fi
% es-noshorthands evita el conflicto de < y > con TikZ
\usepackage[spanish,es-tabla,es-noindentfirst,es-noshorthands]{babel}

% --- geometria y tipografia -----------------------------------------
\usepackage[margin=2.4cm]{geometry}
\usepackage{microtype}
\usepackage{lmodern}
\usepackage{setspace}

% --- matematicas -----------------------------------------------------
\usepackage{amsmath,amssymb,amsthm}
\usepackage{bm}

% --- figuras, tablas, algoritmos ------------------------------------
\usepackage{graphicx}
\usepackage{booktabs}
\usepackage{tabularx}
\newcolumntype{Y}{>{\centering\arraybackslash}X}
\usepackage{xcolor}
\usepackage{tikz}
\usetikzlibrary{arrows.meta,positioning,calc,decorations.pathreplacing}
\usepackage{algorithm}
\usepackage{algpseudocode}
\floatname{algorithm}{Algoritmo}
\renewcommand{\listalgorithmname}{Lista de algoritmos}
\usepackage{caption}
\captionsetup{font=small,labelfont=bf}

% --- referencias cruzadas -------------------------------------------
\usepackage[hidelinks]{hyperref}
\usepackage{cleveref}
\crefname{algorithm}{algoritmo}{algoritmos}
\Crefname{algorithm}{Algoritmo}{Algoritmos}
\crefname{equation}{ecuación}{ecuaciones}
\Crefname{equation}{Ecuación}{Ecuaciones}
\crefname{section}{sección}{secciones}
\Crefname{section}{Sección}{Secciones}
\crefname{figure}{figura}{figuras}
\Crefname{figure}{Figura}{Figuras}
\crefname{table}{tabla}{tablas}
\Crefname{table}{Tabla}{Tablas}
\crefname{proposition}{proposición}{proposiciones}
\Crefname{proposition}{Proposición}{Proposiciones}

% --- entornos de teorema --------------------------------------------
\theoremstyle{plain}
\newtheorem{proposition}{Proposición}
\newtheorem{definition}{Definición}

% --- marca de contenido pendiente -----------------------------------
\newcommand{\todo}[1]{\textcolor{red}{\textbf{[PENDIENTE: #1]}}}
\newcommand{\Sph}{\mathbb{S}}
\newcommand{\Rset}{\mathbb{R}}

\title{\textbf{Separación angular máxima en Transformers de visión:\\
códigos esféricos para una atención diversa y una\\
detección generalizable de medios sintéticos}}

\author{Manuel Muñoz Plá\\
  \small Investigador independiente --- Mairena del Aljarafe, Sevilla, España\\
  \small ORCID: 0009-0000-5714-912X \quad\textbar\quad
  \texttt{mmunozpl@uoc.edu}}

\date{Borrador --- \today}

% =====================================================================
\begin{document}
\maketitle

\begin{abstract}
\noindent
La atención multi-cabeza de los Transformers de visión tiende a
acomodar varias cabezas sobre patrones redundantes, y los detectores de
medios sintéticos pierden exactitud frente a generadores no observados
durante el entrenamiento. Ambos problemas comparten una raíz geométrica:
la disposición de un conjunto de direcciones en el espacio de
representación. Este trabajo propone gobernar esa disposición mediante
\emph{códigos esféricos} ---conjuntos de direcciones unitarias sobre la
hiperesfera con separación angular mínima maximizada--- y deriva dos
contribuciones. La primera regulariza las direcciones representativas de
las cabezas de atención hacia una configuración de máxima separación,
con el objetivo de reducir la redundancia entre cabezas sin alterar la
arquitectura ni añadir parámetros. La segunda fija los prototipos de un
clasificador por generador sobre un código esférico, de modo que el
espacio angular no ocupado denota un margen explícito para el rechazo de
generadores nuevos en régimen de conjunto abierto. El marco se ancla en
el fenómeno de \emph{neural collapse} y emplea el \emph{kissing number}
como caso exacto de validación geométrica. La evaluación, aún
pendiente, contempla ImageNet para la diversidad de atención y el
conjunto GenImage con protocolo \emph{leave-one-generator-out} para la
detección generalizable. \todo{cifras del resumen tras los
experimentos}
\\[4pt]
\noindent\textbf{Palabras clave:} códigos esféricos, separación angular,
atención multi-cabeza, neural collapse, detección de medios sintéticos,
conjunto abierto.
\end{abstract}

% =====================================================================
\section{Introducción}
\label{sec:intro}

El Transformer de visión ha desplazado a la red convolucional como
arquitectura de referencia en clasificación de imagen, y su mecanismo de
atención multi-cabeza se presenta habitualmente como una vía para que el
modelo atienda, en paralelo, a distintas relaciones espaciales. La
evidencia empírica matiza esa lectura: una parte de las cabezas
converge sobre patrones de atención casi indistinguibles, de manera que
la capacidad nominal de la capa no se traduce en diversidad efectiva.
Podar cabezas tras el entrenamiento apenas degrada la exactitud, lo que
denota que la redundancia no es marginal sino estructural.

En un dominio aparentemente alejado, la detección de medios sintéticos
afronta un problema distinto pero igual de persistente. Un detector
entrenado sobre las imágenes de un conjunto de generadores alcanza una
exactitud elevada dentro de ese conjunto y la pierde al enfrentarse a un
generador no observado. El problema abierto del campo no es la
exactitud intradominio, ya saturada, sino la generalización a
arquitecturas generativas nuevas, que aparecen a un ritmo que ningún
conjunto de entrenamiento cerrado puede seguir.

Estos dos problemas, la redundancia de la atención y la fragilidad ante
lo no observado, comparten una raíz que no es evidente a primera vista:
en ambos casos lo que está en juego es \emph{cómo se disponen un
conjunto de direcciones} en un espacio de representación. Las cabezas de
atención pueden caracterizarse por una dirección representativa; las
clases o los generadores, por la dirección del prototipo que el modelo
les asigna. Cuando esas direcciones se amontonan, la capacidad se
desaprovecha y las fronteras de decisión se vuelven frágiles. Cuando se
separan al máximo, cada elemento ocupa una región propia y bien aislada,
y el espacio sobrante queda disponible para lo que el entrenamiento no
preveía.

La geometría que formaliza esa disposición ideal es la de los
\emph{códigos esféricos}: conjuntos de direcciones unitarias sobre la
hiperesfera elegidos de modo que el menor de los ángulos entre pares sea
lo mayor posible. Su caso límite ---el número máximo de direcciones que
admiten una separación angular dada--- es el \emph{kissing number},
un objeto clásico de la geometría de empaquetamientos. Y su contrapartida
en redes neuronales es el fenómeno de \emph{neural collapse}: las
direcciones de clase, sin instrucción explícita, se separan al final del
entrenamiento hasta configurar un marco equiangular. La idea de este
trabajo es no esperar a que la red descubra esa configuración, sino
entregársela fijada de antemano.

\paragraph{Contribuciones.} El trabajo desarrolla dos intervenciones
geométricas sobre el Transformer de visión y un protocolo de evaluación:

\begin{enumerate}
\item Un \textbf{regularizador de atención angularmente diversa}
(\cref{sec:metA}) que conduce las direcciones representativas de las
cabezas hacia una configuración de separación angular máxima, sin
modificar la arquitectura ni incrementar el número de parámetros.
\item Una \textbf{cabeza de prototipos por generador}
(\cref{sec:metB}) cuyas direcciones de clase se fijan sobre un código
esférico, de modo que el espacio angular no asignado actúa como región
de rechazo para generadores no observados, tratando la detección de
medios sintéticos como un problema de conjunto abierto.
\item Un \textbf{protocolo de validación geométrica}
(\cref{sec:metcodes}) que verifica el procedimiento de generación de
códigos contra las configuraciones óptimas conocidas en las dimensiones
$2,3,4,8$ y $24$, y un \textbf{protocolo de evaluación}
(\cref{sec:setup}) con separación \emph{leave-one-generator-out} sobre
el conjunto GenImage.
\end{enumerate}

% =====================================================================
\section{Trabajo relacionado}
\label{sec:related}

\paragraph{Códigos esféricos y empaquetamientos.}
La distribución de puntos sobre la esfera con separación máxima es un
problema clásico, planteado en su forma física por Thomson y en su forma
geométrica por Tammes. El \emph{kissing number} ---el mayor número de
esferas unitarias que pueden tocar a una esfera central sin solaparse---
se conoce de forma exacta solo en las dimensiones $1,2,3,4,8$ y $24$,
estas dos últimas gracias a los retículos $E_8$ y de Leech
\cite{conway1999}. Para dimensiones y cardinalidades arbitrarias se
recurre a códigos casi óptimos obtenidos por minimización de energías de
repulsión \cite{saff1997}.

\paragraph{Neural collapse y clasificadores fijos.}
El fenómeno de \emph{neural collapse} \cite{papyan2020} describe cómo,
en la fase final del entrenamiento, las medias de clase y los vectores
del clasificador convergen a un marco equiangular ajustado. Trabajos
posteriores \cite{zhu2021} lo analizan bajo el modelo de
características no restringidas. Esta regularidad ha motivado el uso de
clasificadores \emph{fijos} con geometría predeterminada
\cite{pernici2021}, que ahorran parámetros y estabilizan el
entrenamiento. El presente trabajo se sitúa en esa línea, pero sustituye
el marco equiangular por un código esférico general y lo extiende a las
proyecciones de atención.

\paragraph{Márgenes angulares.}
SphereFace \cite{liu2017}, CosFace \cite{wang2018} y ArcFace
\cite{deng2019} introdujeron márgenes sobre el ángulo entre la
característica y el prototipo de clase. Estos métodos operan sobre una
geometría angular, pero estiman los prototipos durante el entrenamiento;
aquí los prototipos se fijan con separación garantizada.

\paragraph{Redundancia de la atención multi-cabeza.}
El análisis de la atención multi-cabeza ha mostrado que muchas cabezas
son prescindibles \cite{michel2019} y que cumplen funciones
especializadas e identificables \cite{voita2019}. La diversidad entre
cabezas se ha promovido habitualmente con penalizaciones \emph{ad hoc};
este trabajo la encuadra como un problema de separación angular sobre la
hiperesfera.

\paragraph{Detección de imagen sintética.}
La detección de imagen generada por modelos generativos
\cite{wang2020cnn} y, más recientemente, por modelos de difusión
\cite{corvi2023} se enfrenta a una caída de exactitud al cambiar de
generador. El conjunto GenImage \cite{zhu2023genimage} reúne imágenes de
múltiples generadores y habilita la evaluación de generalización a
arquitecturas no observadas; estrategias como DRCT \cite{chen2024drct}
abordan explícitamente esa transferencia. Ninguna de estas líneas trata
la detección como un problema geométrico de conjunto abierto sobre un
código esférico, que es la posición de este trabajo.

% =====================================================================
\section{Marco teórico}
\label{sec:theory}

\subsection{La hiperesfera: una intuición previa}
\label{sec:theory:intuition}

Cuando varios elementos idénticos que se repelen quedan confinados sobre
una superficie cerrada, no se distribuyen al azar: se acomodan hasta
alcanzar el reposo en que la fuerza neta se anula. Los electrones sobre
una esfera conductora ---el problema que Thomson planteó en 1904--- y las
semillas en el capítulo del girasol obedecen al mismo principio: cada
elemento ocupa la posición que maximiza su distancia angular a los
demás. Ese reposo no es un acto de cómputo, sino el residuo de una
repulsión que se ha dejado actuar hasta el final. De un proceso
enteramente continuo emergen, sin embargo, cantidades enteras ---los doce
vecinos del empaquetamiento tridimensional, los términos de Fibonacci de
la espiral---; el número de elementos no se impone desde fuera, sino que
queda inscrito en la geometría del acomodo. Una red neuronal, al separar
sus clases en el espacio de representación, reproduce ese mismo reposo:
las direcciones asociadas a cada clase se distancian entre sí hasta
configurar una disposición de separación angular máxima, fenómeno que la
literatura reciente denota con el término \emph{neural collapse}. La
diferencia con los sistemas físicos abre una oportunidad: donde la
naturaleza requiere todo el proceso de relajación para alcanzar esa
configuración, un modelo puede recibirla fijada de antemano. Las páginas
siguientes formalizan esa configuración sobre la hiperesfera unitaria
$\Sph^{d-1}$ y precisan la noción de separación angular mínima entre
pares de direcciones.

\begin{figure}[t]
\centering
\begin{tikzpicture}[scale=2.7,>=Stealth]
  \def\R{1}
  % gran circulo de la hiperesfera
  \draw[thick,gray!70] (0,0) circle (\R);
  \draw[dashed,gray!45] (0,0) ellipse ({\R} and 0.34);
  \fill (0,0) circle (0.55pt);
  \node[font=\footnotesize] at (0.12,-0.12) {$O$};
  % direcciones de clase
  \foreach \a in {35,80,150,205,265,315}{
     \draw[->,thick,gray!55] (0,0)--(\a:\R);
     \fill[black] (\a:\R) circle (1.35pt);
  }
  % par con separacion angular minima resaltado
  \draw[->,very thick,blue!70!black] (0,0)--(35:\R);
  \draw[->,very thick,blue!70!black] (0,0)--(80:\R);
  \draw[thick,blue!70!black] (35:0.40) arc (35:80:0.40);
  \node[blue!70!black,font=\footnotesize] at (57:0.60) {$\theta_{\min}$};
  \node[font=\footnotesize] at (35:1.16) {$\bm{u}_i$};
  \node[font=\footnotesize] at (80:1.16) {$\bm{u}_j$};
\end{tikzpicture}
\caption{Esquema bidimensional de un código esférico sobre
$\Sph^{d-1}$. Cada dirección unitaria $\bm{u}_i$ corresponde a una
clase, una cabeza de atención o un generador. La calidad del código se
mide por la separación angular mínima $\theta_{\min}$, el menor de los
ángulos entre pares; un buen código la maximiza.}
\label{fig:sphere}
\end{figure}

\subsection{Códigos esféricos y separación angular}
\label{sec:theory:codes}

\begin{definition}[Código esférico]
Un código esférico de tamaño $K$ en dimensión $d$ es un conjunto
$\mathcal{C}=\{\bm{u}_1,\dots,\bm{u}_K\}\subset\Sph^{d-1}$ de
direcciones unitarias. Su separación angular mínima es
\begin{equation}
\theta_{\min}(\mathcal{C})
  = \min_{i\neq j}\ \arccos\!\big(\langle \bm{u}_i,\bm{u}_j\rangle\big).
\end{equation}
\end{definition}

El objetivo es hallar $\mathcal{C}^\star$ que maximice
$\theta_{\min}$. El mínimo sobre pares no es diferenciable, por lo que en
la práctica se optimiza una energía de repulsión equivalente. Para un
exponente $s>0$, la energía de Riesz
\begin{equation}
E_s(\mathcal{C})
  = \sum_{i\neq j}\ \lVert \bm{u}_i-\bm{u}_j\rVert^{-s}
\label{eq:riesz}
\end{equation}
penaliza la proximidad entre direcciones; sus mínimos aproximan, al
crecer $s$, las configuraciones de separación angular máxima
\cite{saff1997}. La minimización de \cref{eq:riesz} sobre la hiperesfera
admite descenso de gradiente proyectado y constituye el procedimiento de
generación de códigos de \cref{sec:metcodes}.

\subsection{El kissing number como caso exacto}
\label{sec:theory:kissing}

El \emph{kissing number} $\tau(d)$ es el mayor número de esferas
unitarias que pueden tocar simultáneamente a una esfera unitaria central
sin solaparse entre sí. Se conoce de forma exacta solo en dimensiones
contadas: $\tau(2)=6$, $\tau(3)=12$, $\tau(4)=24$, $\tau(8)=240$ y
$\tau(24)=196\,560$ \cite{conway1999}. La conexión con los códigos
esféricos es directa.

\begin{proposition}[Cota de capacidad angular]
\label{prop:kissing}
El número máximo de direcciones unitarias en $\Rset^d$ con separación
angular mínima mayor o igual que $\pi/3$ es exactamente el kissing
number $\tau(d)$.
\end{proposition}

\begin{proof}[Justificación]
Para dos direcciones unitarias,
$\lVert\bm{u}_i-\bm{u}_j\rVert^2 = 2-2\cos\theta_{ij}$, de modo que
$\theta_{ij}\ge\pi/3$ equivale a
$\cos\theta_{ij}\le\tfrac12$ y a
$\lVert\bm{u}_i-\bm{u}_j\rVert\ge 1$. Las bolas de radio $\tfrac12$
centradas en los extremos $\bm{u}_i$ no se solapan y todas son tangentes
a la esfera de radio $\tfrac12$ centrada en el origen; su número máximo
es, por definición, $\tau(d)$ \cite{conway1999}.
\end{proof}

\Cref{prop:kissing} delimita el papel del kissing number en este
trabajo. No es la herramienta operativa ---las dimensiones reales del
modelo no son $8$ ni $24$, ni la cardinalidad de clases coincide con un
$\tau(d)$---, sino el caso exacto que enmarca el objeto general. El
código esférico es lo que la red emplea; el kissing number es la
referencia contra la que se valida el procedimiento que lo genera y la
cota que acota su capacidad.

\subsection{Dos regímenes: holgado y saturado}
\label{sec:theory:regimes}

La separación angular alcanzable por un código de $K$ direcciones en
$\Sph^{d-1}$ depende de la relación entre $K$ y $d$, y conviene
distinguir dos regímenes porque el modelo opera en uno solo de ellos.

\paragraph{Régimen holgado ($K\le d+1$).} Cuando el número de
direcciones no supera $d+1$, existe el \emph{símplex regular}: la
configuración en que todos los pares forman el mismo ángulo, máximo
posible. Sumando las $K$ direcciones unitarias y observando que la
norma al cuadrado es no negativa,
\begin{equation}
0 \le \Big\lVert \textstyle\sum_h \bm{u}_h \Big\rVert^2
  = K + \sum_{h\neq h'}\langle\bm{u}_h,\bm{u}_{h'}\rangle,
\label{eq:simplex}
\end{equation}
se obtiene, para una configuración equiangular de ángulo $\theta$, la
cota $\cos\theta \ge -1/(K-1)$, alcanzada con igualdad por el símplex:
\begin{equation}
\theta_{\mathrm{símplex}}(K)
  = \arccos\!\big(-\tfrac{1}{K-1}\big).
\label{eq:simplexangle}
\end{equation}
Esta separación óptima depende solo de $K$, no de $d$: una vez que la
dimensión alcanza $d=K-1$, añadir más dimensiones deja espacio sin usar
pero no mejora el ángulo. La \cref{eq:simplexangle} es decreciente en
$K$ y está acotada en $(\pi/2,\,\pi]$ ---$180^\circ$ para $K=2$,
$120^\circ$ para $K=3$, y tendiendo a $90^\circ$ por arriba al crecer
$K$, sin alcanzarlo nunca---. El régimen holgado vive, por tanto,
enteramente en ángulos obtusos.

\paragraph{Régimen saturado ($K>d+1$).} Cuando las direcciones superan
$d+1$, el símplex regular deja de existir: no caben tantas direcciones
equiangulares en el espacio disponible. La separación óptima cae por
debajo de $90^\circ$ y la dimensión pasa a ser la restricción activa.
El kissing number es el punto señalado de este régimen ---$K=\tau(d)$
direcciones a exactamente $60^\circ$ (\cref{prop:kissing})---; densidades
mayores producen ángulos aún menores.

\begin{figure}[t]
\centering
\begin{tikzpicture}[scale=1.0,>=Stealth]
  % --- panel izquierdo: regimen holgado, K=3 en 2D, 120 grados ---
  \begin{scope}[shift={(0,0)}]
    \def\R{1.5}
    \draw[gray!55] (0,0) circle (\R);
    \foreach \a in {90,210,330}{
       \draw[->,very thick,blue!65!black] (0,0)--(\a:\R);
       \fill[blue!65!black] (\a:\R) circle (2pt);
    }
    % arco del angulo entre dos direcciones
    \draw[blue!65!black] (90:0.5) arc (90:210:0.5);
    \node[blue!65!black,font=\scriptsize] at (150:0.74) {$120^\circ$};
    \node[font=\scriptsize,align=center] at (0,-2.15)
      {holgado: $K{=}3\le d{+}1$\\símplex, ángulo obtuso};
  \end{scope}
  % --- panel derecho: regimen saturado, K=6 en 2D, 60 grados ---
  \begin{scope}[shift={(5.4,0)}]
    \def\R{1.5}
    \draw[gray!55] (0,0) circle (\R);
    \foreach \a in {0,60,120,180,240,300}{
       \draw[->,very thick,red!70!black] (0,0)--(\a:\R);
       \fill[red!70!black] (\a:\R) circle (2pt);
    }
    \draw[red!70!black] (0:0.5) arc (0:60:0.5);
    \node[red!70!black,font=\scriptsize] at (30:0.74) {$60^\circ$};
    \node[font=\scriptsize,align=center] at (0,-2.15)
      {saturado: $K{=}6=\tau(2)$\\kissing number, ángulo $60^\circ$};
  \end{scope}
\end{tikzpicture}
\caption{Los dos regímenes en dimensión $d{=}2$, donde la
configuración es exacta y representable. Con pocas direcciones para el
espacio disponible (izquierda, $K{=}3$) el símplex las separa al
máximo, en ángulo obtuso. Al saturar el espacio (derecha, $K{=}6$, el
kissing number de $d{=}2$) la separación colapsa a $60^\circ$. El mismo
contraste se da en $d{=}3$ ---tetraedro a $109{,}5^\circ$ frente a
cuboctaedro a $60^\circ$--- y en toda dimensión.}
\label{fig:regimes}
\end{figure}

\paragraph{El modelo opera en régimen holgado.} Las dos intervenciones
de este trabajo se sitúan con holgura en el primer régimen. Las cabezas
de atención de un ViT-B/16 son $H=12$ direcciones y los prototipos de
generador son $K=M+1$ direcciones ---del orden de la decena---, ambos
embebidos en $d=768$. Dado que $12\ll 768$ y $M+1\ll 768$, la dimensión
no es restrictiva. Conviene, sin embargo, distinguir dos variantes del
régimen holgado según la naturaleza de las direcciones.

Los prototipos de generador (\cref{sec:metB}) son direcciones con
orientación definida, y su separación óptima es la del símplex con
signo, $\arccos(-1/(M))\approx 97^\circ$. Las direcciones
representativas de las cabezas (\cref{sec:metA}), en cambio, se definen
como el primer vector singular de $\bm{W}_O^{(h)}$, un objeto
\emph{definido salvo signo}: $\bm{u}_1$ y $-\bm{u}_1$ representan la
misma dirección. Para tales objetos el ángulo con signo no es una
magnitud bien definida, y la separación que sí lo es ---invariante al
signo--- vive en $[0^\circ,90^\circ]$ y vale, para el símplex de $K$
vectores singulares, $\theta^\star_{\mathrm{u}}(K)=\arccos(1/(K-1))$.
Para las $H=12$ cabezas,
$\theta^\star_{\mathrm{u}}=\arccos(1/11)\approx 84{,}8^\circ$ ---el
complementario de los $95{,}2^\circ$ del símplex con signo---. Este es
el valor objetivo $\theta^\star$ del regularizador de \cref{sec:metA}.
El símplex de $12$ vectores no admite, en efecto, ninguna orientación
de signos que deje todos los cosenos mutuos negativos, por lo que el
régimen sin signo no es una elección sino la única descripción
coherente de la separación entre cabezas.

Ninguno de los dos casos toca los $60^\circ$ del kissing number, que
corresponden al régimen saturado y solo intervienen en la validación
geométrica de \cref{sec:metcodes}, donde los códigos esféricos,
generados con orientación definida, sí se miden con signo.

\begin{figure}[t]
\centering
\begin{tikzpicture}[>=Stealth]
  % --- ejes ---
  \def\xa{0}\def\xb{8.4}     % K mapeado: 2..84
  \def\ya{0}\def\yb{4.6}     % theta mapeado: 88..184
  % funciones de mapeo: K en [2,84] -> x ; theta en [88,184] -> y
  % x = (K-2)/82 * 8.4 ; y = (theta-88)/96 * 4.6
  \draw[->] (\xa,\ya) -- (\xb+0.6,\ya)
       node[right,font=\footnotesize] {$K$};
  \draw[->] (\xa,\ya) -- (\xa,\yb+0.3)
       node[above,font=\footnotesize] {$\theta_{\mathrm{símplex}}$};
  % --- asintota en 90 grados ---
  % y(90) = (90-88)/96*4.6 = 0.0958
  \draw[dashed,gray!65] (\xa,0.096) -- (\xb,0.096)
       node[right,font=\scriptsize,gray!55] {$90^\circ$};
  % --- marcas del eje y ---
  \foreach \t/\ty in {90/0.096,120/1.534,150/2.971,180/4.408}{
     \draw (\xa-0.07,\ty)--(\xa+0.07,\ty);
     \node[left,font=\scriptsize] at (\xa-0.1,\ty) {$\t^\circ$};
  }
  % --- marcas del eje x ---
  \foreach \k/\kx in {3/0.102,12/1.024,40/3.892,80/7.700}{
     \draw (\kx,\ya-0.07)--(\kx,\ya+0.07);
     \node[below,font=\scriptsize] at (\kx,\ya-0.1) {\k};
  }
  % --- curva theta_simplex(K), puntos calculados ---
  \draw[very thick,blue!65!black] plot[smooth] coordinates {
    (0.000,4.408) (0.102,1.534) (0.205,1.030) (0.307,0.790)
    (0.410,0.649) (0.614,0.490) (1.024,0.346) (1.844,0.241)
    (3.892,0.166) (7.700,0.132) };
  % --- punto K=12 sobre la curva con signo: 95,2 grados ---
  \fill[blue!65!black] (1.024,0.346) circle (2.4pt);
  \draw[blue!65!black,thin] (1.024,0.346)
       -- ++(0.7,0.7)
       node[right,font=\scriptsize,align=left]
       {$K{=}12$: símplex con signo\\$\approx\!95{,}2^\circ$};
\end{tikzpicture}
\caption{Separación angular óptima del símplex con signo
$\theta_{\mathrm{símplex}}(K)=\arccos(-1/(K-1))$ en función del número
de direcciones $K$, válida en régimen holgado ($K\le d+1$). La curva
decrece de forma monótona desde $180^\circ$ y se aproxima
asintóticamente a $90^\circ$ sin alcanzarlo: en este régimen la
separación óptima con signo es siempre obtusa. El punto azul marca el
caso $K{=}12$ ($\approx95{,}2^\circ$), aplicable a los prototipos de
generador y a la validación de códigos. Las direcciones
representativas de las cabezas de atención, en cambio, son vectores
singulares sin signo: su separación óptima es la reflexión respecto a
$90^\circ$ de la curva, $\arccos(1/(K-1))$, que para $K{=}12$ vale
$\approx84{,}8^\circ$.}
\label{fig:simplexcurve}
\end{figure}

\subsection{Conexión con neural collapse}
\label{sec:theory:nc}

El fenómeno de \emph{neural collapse} \cite{papyan2020} establece que,
al final del entrenamiento, los vectores del clasificador de un modelo
bien entrenado convergen a un marco equiangular ajustado: una
configuración de separación angular máxima para $K\le d+1$ clases. La
red, por tanto, tiende espontáneamente hacia el código esférico. Las dos
contribuciones de este trabajo parten de esa observación y la
anticipan: en lugar de esperar a que la geometría emerja del proceso de
relajación, se fija o se regulariza desde el inicio. En el régimen
holgado en que opera el modelo (\cref{sec:theory:regimes}) esa
disposición óptima es la del símplex regular; el marco equiangular del
\emph{neural collapse} y el código esférico de Riesz coinciden ahí
numéricamente, observación que la \cref{sec:setup} retoma al describir
las ablaciones.

% =====================================================================
\section{Método}
\label{sec:method}

\begin{figure}[t]
\centering
\begin{tikzpicture}[
  >=Stealth,
  block/.style={draw,rounded corners,align=center,
    minimum height=11mm,minimum width=17mm,font=\footnotesize},
  every node/.style={font=\footnotesize}]
  \node[block] (img) {Imagen};
  \node[block,right=8mm of img] (patch) {Patch\\embedding};
  \node[block,right=8mm of patch,fill=blue!6] (enc)
       {Encoder ViT\\$L$ bloques};
  \node[block,right=8mm of enc] (cls) {Embedding\\\texttt{[CLS]}\;$\bm{z}$};
  \node[block,right=8mm of cls,fill=blue!6] (head)
       {Cabeza de\\prototipos};
  \node[block,right=8mm of head] (out) {Logits\\+ rechazo};
  \foreach \a/\b in {img/patch,patch/enc,enc/cls,cls/head,head/out}
     {\draw[->] (\a)--(\b);}
  % llamada A: atencion angularmente diversa
  \node[block,below=11mm of enc,fill=blue!6,minimum width=30mm]
       (mha) {Contribución A\\atención angularmente diversa};
  \draw[->,dashed] (mha)--(enc);
  % llamada B: prototipos sobre codigo esferico
  \node[block,below=11mm of head,fill=blue!6,minimum width=30mm]
       (proto) {Contribución B\\prototipos en código esférico};
  \draw[->,dashed] (proto)--(head);
\end{tikzpicture}
\caption{Visión general. El Transformer de visión recibe dos
intervenciones geométricas independientes: la contribución A regulariza
las direcciones de las cabezas de atención dentro del encoder; la
contribución B fija los prototipos de la cabeza de clasificación sobre
un código esférico y habilita el rechazo de generadores no observados.}
\label{fig:overview}
\end{figure}

\subsection{Generación y validación de códigos esféricos}
\label{sec:metcodes}

Dado un par $(d,K)$, el código $\mathcal{C}$ se obtiene minimizando la
energía de Riesz de \cref{eq:riesz} mediante descenso de gradiente
proyectado sobre $\Sph^{d-1}$: en cada paso el gradiente se proyecta al
espacio tangente a la hiperesfera, se aplica un paso del optimizador
Adam y las direcciones se renormalizan a norma unidad.\footnote{El
optimizador adaptativo y la proyección tangente mejoran la convergencia
frente al descenso estocástico ingenuo en $d\gg 1$, donde la energía de
Riesz presenta un paisaje altamente no convexo: los gradientes se
disparan con pares próximos en la inicialización aleatoria y, sin
proyección tangente, la dinámica deriva por la componente radial.} El
procedimiento se resume en \cref{alg:codes}.

\begin{algorithm}[t]
\caption{Generación de un código esférico por descenso proyectado}
\label{alg:codes}
\begin{algorithmic}[1]
\Require dimensión $d$, tamaño $K$, exponente $s$, pasos $T$, tasa
inicial $\eta_0$
\State inicializar $\bm{u}_1,\dots,\bm{u}_K$ aleatorias en $\Sph^{d-1}$
\State inicializar el optimizador Adam con tasa $\eta_0$
\State inicializar el planificador \emph{cosine annealing}
\For{$t = 1$ \textbf{hasta} $T$}
  \State calcular $\bm{g}_i = \nabla_{\bm{u}_i} E_s$
         \Comment{energía de repulsión}
  \State $\bm{g}_i^{\mathrm{tan}}
         \gets \bm{g}_i - \langle\bm{g}_i,\bm{u}_i\rangle\,\bm{u}_i$
         \Comment{proyección tangente}
  \State $\bm{u}_i \gets$ paso de Adam sobre $\bm{u}_i$
         con $\bm{g}_i^{\mathrm{tan}}$
  \State $\bm{u}_i \gets \bm{u}_i / \lVert\bm{u}_i\rVert$
         \Comment{reproyección a la hiperesfera}
  \State actualizar la tasa con el planificador
\EndFor
\State \Return $\mathcal{C}=\{\bm{u}_1,\dots,\bm{u}_K\}$
\end{algorithmic}
\end{algorithm}

\paragraph{Validación geométrica.} Antes de su uso en la red, el
procedimiento se verifica en las dimensiones de kissing number conocido.
Para $(d,K)\in\{(2,6),(3,12),(4,24),(8,240),(24,196560)\}$ se comprueba
que el código recuperado satisface la cota del kissing number,
$\theta_{\min}\ge\pi/3$, dentro de una tolerancia numérica. Para las
dimensiones reales del modelo, sin solución exacta, se reporta la brecha
respecto a las cotas superiores conocidas. Este control de sanidad
precede a todo experimento con redes y no consume cómputo de GPU.

El descenso desde inicialización aleatoria no recupera, sin embargo,
todas las configuraciones óptimas. Para $(d,K){=}(4,24)$ se estanca de
forma reproducible ---verificado sobre seis semillas, cinco valores del
exponente $s$, ruido térmico recocido e inicialización ortogonal--- en
un mínimo local con $\theta_{\min}\approx55{,}23^\circ$ y energía a un
$0{,}007$ por ciento del óptimo de $D_4$: la configuración $D_4$ existe, pero
su cuenca de atracción desde inicialización aleatoria es demasiado
estrecha. Para $(d,K){=}(3,12)$ el descenso converge al icosaedro
regular, con $\theta_{\min}\approx63{,}43^\circ$, en lugar del
cuboctaedro a $60^\circ$; el icosaedro es más estable bajo la energía
de Riesz y satisface igualmente la cota del kissing
($63{,}43^\circ\ge60^\circ$). Por ello, para los pares con
configuración óptima cerrada conocida ---$D_4$ en $d{=}4$, $E_8$ en
$d{=}8$, raíces del retículo de Leech en $d{=}24$--- el protocolo
arranca desde la configuración canónica y verifica que el descenso la
conserva; para las dimensiones reales del modelo ---$d{=}768$,
$K\in\{9,12\}$--- la inicialización es aleatoria y la separación
recuperada se reporta en la \cref{tab:geo}.

\subsection{Contribución A: atención angularmente diversa}
\label{sec:metA}

Una capa de atención multi-cabeza con $H$ cabezas dispone, para cada
cabeza $h$, de una proyección de salida
$\bm{W}_O^{(h)}\in\Rset^{d\times d_h}$. Se define la \emph{dirección
representativa} de la cabeza $h$ como el primer vector singular
izquierdo de $\bm{W}_O^{(h)}$:
\begin{equation}
\bm{r}_h = \bm{u}_1\!\big(\bm{W}_O^{(h)}\big),
\qquad
\bm{u}_1 \;\text{primer vector singular izquierdo},
\label{eq:headdir}
\end{equation}
es decir, la dirección del subespacio de salida en que la cabeza
deposita más energía. Esta elección no es arbitraria. Un estudio
preliminar sobre el modelo afinado contrasta dos resúmenes posibles de
la cabeza ---el primer vector singular y la media normalizada de las
columnas de $\bm{W}_O^{(h)}$--- frente a una medida directa de
diversidad funcional: la similitud entre los mapas de atención de cada
par de cabezas. El primer vector singular correla con esa similitud
funcional de forma apreciable (coeficiente agregado $\approx0{,}35$
sobre las doce capas), mientras que la media de columnas resulta
prácticamente ortogonal al comportamiento de la cabeza (coeficiente
$\approx0{,}02$).\footnote{El primer vector singular está definido
salvo signo; la comparación funcional y el regularizador de
\cref{eq:rdiv} emplean por ello el valor absoluto del producto
escalar. La correlación, aunque clara en agregado, es parcial: la
similitud entre mapas de atención está dominada por estructura común
---atención al token de clase y a los parches vecinos--- y el resumen
no captura la totalidad de la diversidad funcional.} La media de
columnas mide la rotación del centroide de $\bm{W}_O^{(h)}$, una
cantidad que un modelo puede alterar sin modificar sus mapas de
atención; el primer vector singular capta, en cambio, el subespacio
dominante de la proyección.

El conjunto $\{\bm{r}_1,\dots,\bm{r}_H\}$ es un código esférico en
$\Rset^d$. Dado que cada $\bm{r}_h$ es un vector singular, definido
salvo signo, la intervención añade a la función de pérdida un término
que empuja ese conjunto hacia la separación angular máxima usando el
valor absoluto del producto escalar:
\begin{equation}
\mathcal{R}_{\mathrm{div}}
  = \sum_{h\neq h'}
    \max\!\big(0,\;
      \lvert\langle\bm{r}_h,\bm{r}_{h'}\rangle\rvert
      - \cos\theta^\star\big),
\label{eq:rdiv}
\end{equation}
donde $\theta^\star=\arccos(1/(H-1))\approx 84{,}8^\circ$ es la
separación óptima sin signo del símplex de $H$ vectores singulares
---el régimen sin signo descrito en \cref{sec:theory:regimes}---. El
valor absoluto es necesario, no opcional: dos cabezas con vectores
singulares casi antiparalelos son funcionalmente equivalentes, y el
producto escalar con signo las leería erróneamente como bien
separadas. Nótese que $\cos\theta^\star=1/(H-1)>0$, de modo que la
penalización se anula exactamente cuando las cabezas alcanzan el
símplex sin signo; no es un umbral estructural. El término penaliza
únicamente los pares por debajo del objetivo, no impone ninguna
arquitectura nueva y no añade parámetros entrenables. Se contempla
además una variante \emph{dura}, en la que la proyección
$\bm{W}_O^{(h)}$ se somete tras cada paso a una rotación ortogonal que
alinea su primer vector singular con la dirección asignada del código;
la comparación entre la variante blanda de \cref{eq:rdiv} y la dura
forma parte de la ablación de \cref{sec:setup}.

\subsection{Contribución B: prototipos de generador en conjunto abierto}
\label{sec:metB}

Para la detección de medios sintéticos, el espacio de etiquetas de
entrenamiento es $\{\text{real}\}\cup\{g_1,\dots,g_M\}$, con $M$
generadores observados. La cabeza de clasificación asigna a cada
etiqueta un prototipo unitario, y el conjunto de los $M+1$ prototipos se
fija sobre un código esférico $\mathcal{C}$ generado por
\cref{alg:codes}. Dado el embedding \texttt{[CLS]} normalizado
$\hat{\bm{z}}=\bm{z}/\lVert\bm{z}\rVert$, la puntuación de la clase $k$
es el coseno $s_k=\langle\hat{\bm{z}},\bm{p}_k\rangle$, y el
entrenamiento emplea una pérdida de margen angular que atrae cada
embedding hacia el prototipo de su etiqueta.

\paragraph{Régimen de conjunto abierto.} La separación angular máxima
entre prototipos deja, entre las celdas de decisión, regiones angulares
no asignadas a ninguna etiqueta. Una imagen procedente de un generador
no observado tiende a proyectarse en esas regiones en lugar de invadir
la celda de la clase \emph{real}. Se definen dos decisiones:

\begin{itemize}
\item \textbf{Real frente a sintético:} se compara la puntuación del
prototipo \emph{real} con la mayor puntuación entre los prototipos de
generador; un margen a favor de los segundos denota imagen sintética.
\item \textbf{Generador novel:} si $\max_k s_k < \tau$ para un umbral
$\tau$, la imagen se marca como procedente de un generador no
catalogado, sin perder la decisión binaria anterior.
\end{itemize}

\begin{figure}[t]
\centering
\begin{tikzpicture}[scale=2.7,>=Stealth]
  \def\R{1}
  \draw[thick,gray!70] (0,0) circle (\R);
  \fill (0,0) circle (0.55pt);
  % prototipos conocidos: real + 4 generadores, separacion maxima
  \foreach \a/\lab in {90/real,162/$g_1$,234/$g_2$,306/$g_3$,18/$g_4$}{
     % cono de aceptacion (semiangulo 26 grados)
     \fill[blue!10]
        (0,0) -- (\a-26:\R) arc (\a-26:\a+26:\R) -- cycle;
     \draw[->,thick,blue!65!black] (0,0)--(\a:\R);
     \fill[blue!65!black] (\a:\R) circle (1.5pt);
     \node[font=\scriptsize,blue!65!black] at (\a:1.18) {\lab};
  }
  % muestra de generador no observado, cae en region no asignada
  \draw[->,very thick,red!75!black] (0,0)--(126:\R);
  \fill[red!75!black] (126:\R) circle (1.5pt);
  \node[font=\scriptsize,red!75!black,align=center]
       at (126:1.30) {generador\\no observado};
\end{tikzpicture}
\caption{Geometría del rechazo en conjunto abierto. Los prototipos de
las clases conocidas ---\emph{real} y generadores $g_1\dots g_4$---
ocupan direcciones de separación angular máxima; cada uno define un cono
de aceptación (sombreado). Una imagen de un generador no observado tiende
a proyectarse en el espacio angular no asignado y queda fuera de todos
los conos, lo que habilita su rechazo sin contaminar la celda
\emph{real}.}
\label{fig:openset}
\end{figure}

\subsection{Objetivo de entrenamiento}
\label{sec:loss}

La función de pérdida combina la tarea, el regularizador de diversidad y
la clasificación por prototipos:
\begin{equation}
\mathcal{L}
  = \mathcal{L}_{\mathrm{tarea}}
  + \lambda_A\,\mathcal{R}_{\mathrm{div}}
  + \lambda_B\,\mathcal{L}_{\mathrm{proto}},
\label{eq:loss}
\end{equation}
con pesos $\lambda_A,\lambda_B\ge 0$. En los experimentos de diversidad
de atención sobre clasificación de imagen, $\lambda_B=0$ y
$\mathcal{L}_{\mathrm{tarea}}$ es la entropía cruzada habitual. En los
experimentos de detección, $\mathcal{L}_{\mathrm{proto}}$ asume el papel
de clasificador y $\mathcal{R}_{\mathrm{div}}$ se mantiene como término
opcional cuya aportación se cuantifica por ablación.

% =====================================================================
\section{Configuración experimental}
\label{sec:setup}

Esta sección fija el protocolo; los resultados se incorporarán una vez
ejecutados los experimentos.

\paragraph{Conjuntos de datos.}
Para la contribución A se emplea ImageNet-100, el subconjunto de 100
clases de ImageNet. Cien clases aportan riqueza angular suficiente para
que la diversidad entre cabezas sea medible ---mientras que diez la
dejarían trivial---, y el fenómeno que la contribución A estudia es
geométrico: depende de la relación entre el número de cabezas y la
dimensión del modelo, no del número de clases, por lo que ImageNet-100
es representativo sin necesidad de la escala completa. Para la
contribución B se emplea GenImage \cite{zhu2023genimage}, que reúne
imágenes reales y sintéticas de múltiples generadores, con el protocolo
\emph{leave-one-generator-out} descrito más abajo; DRCT
\cite{chen2024drct} se reserva como conjunto de prueba externo, solo en
inferencia, para medir transferencia.

\paragraph{Protocolo \emph{leave-one-generator-out}.}
Para cada generador $g$ de GenImage se entrena un detector con las
imágenes de todos los generadores excepto $g$, y se evalúa la detección
sobre $g$. La media sobre los $M$ generadores excluidos cuantifica la
generalización a arquitecturas no observadas, que es el objeto central
de la contribución B.

\paragraph{Arquitectura y entorno.}
La columna vertebral es un Transformer de visión ViT-B/16. La
contribución A parte de pesos preentrenados de ViT-B/16 (\texttt{timm})
y se afina sobre ImageNet con la pérdida $\mathcal{L}$ de la
\cref{eq:loss}; el regularizador $\mathcal{R}_{\mathrm{div}}$ actúa así
como un aditivo geométrico sobre una columna ya competente, lo que
aísla su efecto mejor que un entrenamiento conjunto desde cero ---que
con el cómputo disponible no alcanzaría una exactitud competitiva sin
la receta completa de entrenamiento prolongado---. La contribución B
parte igualmente de pesos preentrenados y se afina sobre GenImage a una
resolución de 256\,px, suficiente para preservar los artefactos de alta
frecuencia de los generadores de difusión. El entorno de cómputo es una
GPU NVIDIA RTX~5090 de 32\,GB, que admite el ViT-B en una sola tarjeta
sin recurrir a \emph{gradient checkpointing}.

\paragraph{Métricas.}
La diversidad de atención se mide por la redundancia entre cabezas
---similitud media entre las direcciones representativas--- y por la
entropía de los mapas de atención, además de la exactitud top-1. La
detección se evalúa por exactitud binaria, AUROC y, en régimen de
conjunto abierto, por la curva de rechazo de clasificación abierta
(OSCR) y el AUROC de separación entre clases conocidas y generadores no
observados. La validación geométrica reporta la separación angular
alcanzada frente al óptimo conocido.

\paragraph{Ablaciones previstas.}
(i) Variante blanda (\cref{eq:rdiv}) frente a variante dura del código
de atención. (ii) Prototipos fijos frente a prototipos inicializados con
el código y posteriormente afinados. (iii) Código esférico frente a
inicialización aleatoria; en las dimensiones del modelo ($K\ll d$) el
código de Riesz coincide numéricamente con el marco equiangular
ajustado ---ambos con $\theta_{\min}=\arccos(-1/(K-1))$---, por lo que
esa tercera configuración no se reporta como condición separada. (iv)
Sensibilidad a los pesos $\lambda_A,\lambda_B$ y al exponente $s$ de la
energía de Riesz.

% =====================================================================
\section{Resultados}
\label{sec:results}

Se reportan a continuación los resultados de la contribución~A sobre
ImageNet-100. Los de la validación geométrica (\cref{tab:geo}) y la
contribución~B (\cref{tab:detect}) se incorporarán tras su ejecución.

\subsection{Contribución A: diversidad de atención}
\label{sec:results:A}

Las tres configuraciones ---base sin regularizador, blanda con
$\lambda_A=10^{-3}$ y dura por reproyección--- se afinan 30 épocas sobre
ImageNet-100 desde pesos preentrenados de ViT-B/16. La \cref{tab:divattn}
resume el estado final y la \cref{fig:thetaevol} la evolución por época.

\begin{table}[h]
\centering
\caption{Contribución A sobre ImageNet-100, estado final tras 30 épocas.
$\theta_{\min}$ es la separación angular mínima sin signo entre cabezas;
$\bar{s}$ es la similitud funcional media entre mapas de atención.}
\label{tab:divattn}
\begin{tabularx}{0.96\linewidth}{X Y Y Y Y}
\toprule
Configuración & $\theta_{\min}$ $\uparrow$ &
Redundancia $\downarrow$ & $\bar{s}$ func. $\downarrow$ &
Top-1 $\uparrow$ \\
\midrule
Base ($\lambda_A=0$)                & $52{,}0^\circ$ & $0{,}175$ &
$0{,}693$ & $0{,}925$ \\
Blanda ($\lambda_A=10^{-3}$)        & $85{,}0^\circ$ & $0{,}034$ &
$0{,}688$ & $\mathbf{0{,}926}$ \\
Dura (reproyección)                 & $84{,}8^\circ$ & $0{,}091$ &
--- & $0{,}877$ \\
\bottomrule
\end{tabularx}
\end{table}

\begin{figure}[t]
\centering
\begin{tikzpicture}[>=Stealth]
  \def\xb{10}\def\yb{5}
  \draw[->] (0,0)--(\xb+0.4,0) node[right,font=\footnotesize]{época};
  \draw[->] (0,0)--(0,\yb+0.4)
       node[above,font=\footnotesize]{$\theta_{\min}$};
  % marcas y: 45,60,75,90 -> (t-45)/45*5
  \foreach \t/\ty in {45/0,60/1.667,75/3.333,90/5.0}{
     \draw (-0.07,\ty)--(0.07,\ty);
     \node[left,font=\scriptsize] at (-0.1,\ty) {$\t^\circ$};}
  % marcas x: 1,10,20,30 -> (e-1)/29*10
  \foreach \e/\ex in {1/0,10/3.103,20/6.552,30/10.0}{
     \draw (\ex,-0.07)--(\ex,0.07);
     \node[below,font=\scriptsize] at (\ex,-0.1) {\e};}
  % objetivo sin signo 84,8 -> 4.422
  \draw[dashed,gray!55] (0,4.422)--(\xb,4.422)
       node[right,font=\scriptsize,gray!55]{$\theta^\star$};
  % dura (clavada)
  \draw[thick,red!70!black]
    (0,4.422)--(10,4.422);
  % blando
  \draw[very thick,blue!65!black] plot[smooth] coordinates {
    (0.000,2.900)(0.345,3.156)(0.690,3.389)(1.034,3.722)(1.379,4.122)
    (1.724,4.411)(2.069,4.433)(2.414,4.444)(3.103,4.444)(3.793,4.444)
    (4.483,4.433)(5.172,4.444)(5.862,4.456)(6.552,4.444)(7.241,4.444)
    (7.931,4.444)(8.621,4.444)(9.310,4.444)(10.000,4.444)};
  % base
  \draw[thick,gray!60] plot[smooth] coordinates {
    (0.000,0.389)(0.345,0.511)(0.690,0.633)(1.379,0.644)(2.069,0.711)
    (2.759,0.822)(3.448,0.844)(4.138,0.833)(4.828,0.800)(5.517,0.822)
    (6.207,0.789)(6.897,0.778)(7.586,0.756)(8.276,0.778)(8.966,0.767)
    (10.000,0.767)};
  \node[font=\scriptsize,blue!65!black] at (8.3,4.72) {blanda};
  \node[font=\scriptsize,red!70!black] at (8.3,4.16) {dura};
  \node[font=\scriptsize,gray!60] at (8.3,1.02) {base};
\end{tikzpicture}
\caption{Evolución de la separación angular mínima sin signo entre
cabezas a lo largo del afinado. La base, sin regularizador, se estanca
en torno a $52^\circ$ ---el techo del \emph{neural collapse}
espontáneo---. La blanda alcanza el óptimo $\theta^\star\approx84{,}8^\circ$
hacia la época~6 y lo mantiene. La dura lo impone desde la primera
época por reproyección.}
\label{fig:thetaevol}
\end{figure}

Tres observaciones se desprenden de los datos.

Primero, el afinado por sí solo separa las cabezas, pero solo
parcialmente: la base alcanza $\theta_{\min}\approx52^\circ$ y se
estanca, muy por debajo del óptimo de $84{,}8^\circ$. Es el
\emph{neural collapse} esperado, con un techo que deja redundancia
angular sin eliminar. Esto justifica que la intervención de la
contribución~A no sea redundante con el fenómeno espontáneo: el
\emph{collapse} no llega al óptimo.

Segundo, el regularizador blando alcanza el óptimo del símplex sin
signo y reduce la redundancia a la mitad ---de $0{,}175$ a
$0{,}034$--- sin coste en exactitud: el \emph{top-1} de validación
($0{,}926$) iguala al de la base ($0{,}925$). La separación angular
máxima entre cabezas es, por tanto, alcanzable sin penalización de
rendimiento. La variante dura llega a la misma geometría pero degrada
la exactitud casi cinco puntos ($0{,}877$), lo que denota que
incentivar la separación mediante un término suave es preferible a
imponerla por reproyección: el camino hacia la configuración óptima
importa tanto como la configuración.

Tercero, y con la cautela debida, el efecto sobre la diversidad
\emph{funcional} de la atención es marginal. La similitud media entre
mapas de atención apenas desciende ---de $0{,}693$ a $0{,}688$--- y la
reducción se concentra casi enteramente en las capas tardías
($-0{,}035$ en la última capa), mientras varias capas intermedias no
mejoran o empeoran levemente. La separación angular de las direcciones
de salida $\bm{W}_O$ y la diversidad de comportamiento de la atención
resultan, así, en gran medida \emph{desacopladas}: reorganizar la
geometría de $\bm{W}_O$ hacia la máxima separación apenas altera a qué
atiende cada cabeza. Este desacople es en sí mismo un hallazgo, y
orienta la continuación natural del trabajo (\cref{sec:discussion}).

\subsection{Validación geométrica y contribución B}
\label{sec:results:resto}


\begin{table}[h]
\centering
\caption{Validación geométrica: separación angular recuperada frente a
la cota del kissing number.}
\label{tab:geo}
\begin{tabularx}{0.86\linewidth}{l c c c Y}
\toprule
Dimensión $d$ & Tamaño $K$ & Cota kissing ($\theta_{\min}\!\ge$) &
$\theta_{\min}$ recuperado & Conforme \\
\midrule
2  & 6      & $60^\circ$ & \todo{} & \todo{} \\
3  & 12     & $60^\circ$ & \todo{} & \todo{} \\
4  & 24     & $60^\circ$ & \todo{} & \todo{} \\
8  & 240    & $60^\circ$ & \todo{} & \todo{} \\
24 & 196560 & $60^\circ$ & \todo{} & \todo{} \\
\bottomrule
\end{tabularx}
\end{table}

\begin{table}[h]
\centering
\caption{Contribución B. Detección \emph{leave-one-generator-out} sobre
GenImage y transferencia a DRCT.}
\label{tab:detect}
\begin{tabularx}{0.96\linewidth}{X Y Y Y}
\toprule
Configuración & AUROC GenImage (LOGO) $\uparrow$ &
OSCR $\uparrow$ & AUROC transferencia DRCT $\uparrow$ \\
\midrule
Cabeza lineal base                 & \todo{} & \todo{} & \todo{} \\
+ prototipos en código esférico    & \todo{} & \todo{} & \todo{} \\
+ rechazo de conjunto abierto      & \todo{} & \todo{} & \todo{} \\
\bottomrule
\end{tabularx}
\end{table}

\todo{figura de la validación geométrica ---separación alcanzada por
dimensión--- y figura de las curvas OSCR por generador excluido.}

% =====================================================================
\section{Discusión y limitaciones}
\label{sec:discussion}

El marco descansa sobre una distinción que conviene no difuminar. Las
configuraciones de separación angular máxima solo se conocen de forma
exacta en dimensiones contadas; en las dimensiones reales del modelo se
trabaja con códigos casi óptimos, y la brecha respecto al óptimo
teórico ---cuantificada en \cref{tab:geo}--- denota el margen de
imprecisión que el método arrastra. El kissing number es el caso ideal
de referencia, no la herramienta universal.

Una segunda cuestión abierta es si las direcciones del código deben
fijarse e inmovilizarse o solo inicializar un ajuste posterior. La
primera opción ofrece el argumento teórico más limpio y el ahorro de
parámetros; la segunda suele rendir algo más en exactitud. El trabajo no
prejuzga la respuesta y la traslada a la ablación de \cref{sec:setup}.

Por último, la dirección representativa de una cabeza de atención
---definida en \cref{eq:headdir} como primer vector singular de
$\bm{W}_O^{(h)}$--- es un resumen entre varios posibles. La elección se
apoya en el estudio preliminar de \cref{sec:metA}, que la prefiere
sobre la media de columnas por su correlación apreciablemente mayor con
la similitud entre mapas de atención. Esa correlación es, sin embargo,
parcial: el resumen captura una fracción de la diversidad funcional, no
su totalidad, y su ventaja sobre la media de columnas no es uniforme
entre capas. La cuantía con que la regularización de un resumen así
definido se traduce en diversidad efectiva de la atención es, por
tanto, una cuestión que la evaluación debe medir, no un hecho que el
marco garantice.

Los resultados de la \cref{sec:results:A} resuelven esa cuestión en
sentido matizado: la regularización alcanza el óptimo geométrico sin
coste de exactitud, pero su efecto sobre la diversidad funcional de la
atención es marginal. La interpretación más plausible es que la función
de una cabeza ---a qué atiende--- no reside en la proyección de salida
$\bm{W}_O$, que solo proyecta de vuelta el resultado, sino en el
mecanismo que decide la atención, $\mathrm{softmax}(QK^\top)$. Separar
las direcciones de $\bm{W}_O$ reorganiza la salida sin tocar ese
mecanismo, de donde el desacople observado. Inducir separación sobre el
subespacio $Q\cdot K$ ---y el estudio de atribución por componente que
lo precede--- es la continuación natural de este trabajo; a diferencia
de $\bm{W}_O$, ese subespacio define una forma bilineal y no un conjunto
de direcciones, por lo que exige extender el marco más allá de los
códigos esféricos.

% =====================================================================
\section{Conclusión}
\label{sec:conclusion}

El trabajo encuadra dos problemas en apariencia distantes ---la
redundancia de la atención multi-cabeza y la fragilidad de los
detectores de medios sintéticos frente a generadores no observados---
como un único problema geométrico: la disposición de un conjunto de
direcciones sobre la hiperesfera. Sobre esa base propone un regularizador
de atención angularmente diversa y una cabeza de prototipos por
generador anclada en un código esférico, ambos sin coste arquitectónico
y con el espacio angular no asignado como margen explícito para lo no
observado. La validación geométrica frente al kissing number y la
evaluación con protocolo \emph{leave-one-generator-out} quedan
especificadas; su ejecución determinará el alcance empírico de la
propuesta.

% =====================================================================
\section*{Disponibilidad de código y datos}
\addcontentsline{toc}{section}{Disponibilidad de código y datos}
El código de generación de códigos esféricos, de las dos intervenciones
y de los protocolos de evaluación se publicará en un repositorio abierto
junto con la versión revisada de este preprint. \todo{enlace al
repositorio y DOI del archivo de pesos}

% =====================================================================
\begin{thebibliography}{99}

\bibitem{conway1999}
J.~H. Conway y N.~J.~A. Sloane.
\emph{Sphere Packings, Lattices and Groups}.
Springer, 3.ª ed., 1999.

\bibitem{saff1997}
E.~B. Saff y A.~B.~J. Kuijlaars.
Distributing many points on a sphere.
\emph{The Mathematical Intelligencer}, 19(1):5--11, 1997.

\bibitem{papyan2020}
V.~Papyan, X.~Y. Han y D.~L. Donoho.
Prevalence of neural collapse during the terminal phase of deep
learning training.
\emph{Proceedings of the National Academy of Sciences},
117(40):24652--24663, 2020.

\bibitem{zhu2021}
Z.~Zhu, T.~Ding, J.~Zhou, X.~Li, C.~You, J.~Sulam y Q.~Qu.
A geometric analysis of neural collapse with unconstrained features.
\emph{Advances in Neural Information Processing Systems} (NeurIPS), 2021.

\bibitem{pernici2021}
F.~Pernici, M.~Bruni, C.~Baecchi y A.~Del~Bimbo.
Regular polytope networks.
\emph{IEEE Transactions on Neural Networks and Learning Systems}, 2021.

\bibitem{liu2017}
W.~Liu, Y.~Wen, Z.~Yu, M.~Li, B.~Raj y L.~Song.
SphereFace: deep hypersphere embedding for face recognition.
\emph{IEEE Conference on Computer Vision and Pattern Recognition}
(CVPR), 2017.

\bibitem{wang2018}
H.~Wang, Y.~Wang, Z.~Zhou, X.~Ji, D.~Gong, J.~Zhou, Z.~Li y W.~Liu.
CosFace: large margin cosine loss for deep face recognition.
\emph{IEEE Conference on Computer Vision and Pattern Recognition}
(CVPR), 2018.

\bibitem{deng2019}
J.~Deng, J.~Guo, N.~Xue y S.~Zafeiriou.
ArcFace: additive angular margin loss for deep face recognition.
\emph{IEEE Conference on Computer Vision and Pattern Recognition}
(CVPR), 2019.

\bibitem{michel2019}
P.~Michel, O.~Levy y G.~Neubig.
Are sixteen heads really better than one?
\emph{Advances in Neural Information Processing Systems} (NeurIPS), 2019.

\bibitem{voita2019}
E.~Voita, D.~Talbot, F.~Moiseev, R.~Sennrich y I.~Titov.
Analyzing multi-head self-attention: specialized heads do the heavy
lifting, the rest can be pruned.
\emph{Annual Meeting of the Association for Computational Linguistics}
(ACL), 2019.

\bibitem{dosovitskiy2021}
A.~Dosovitskiy, L.~Beyer, A.~Kolesnikov \emph{et al.}
An image is worth 16x16 words: Transformers for image recognition at
scale.
\emph{International Conference on Learning Representations} (ICLR), 2021.

\bibitem{vaswani2017}
A.~Vaswani, N.~Shazeer, N.~Parmar \emph{et al.}
Attention is all you need.
\emph{Advances in Neural Information Processing Systems} (NeurIPS), 2017.

\bibitem{wang2020cnn}
S.-Y. Wang, O.~Wang, R.~Zhang, A.~Owens y A.~A. Efros.
CNN-generated images are surprisingly easy to spot\dots\ for now.
\emph{IEEE Conference on Computer Vision and Pattern Recognition}
(CVPR), 2020.

\bibitem{corvi2023}
R.~Corvi, D.~Cozzolino, G.~Zingarini, G.~Poggi, K.~Nagano y
L.~Verdoliva.
On the detection of synthetic images generated by diffusion models.
\emph{IEEE International Conference on Acoustics, Speech and Signal
Processing} (ICASSP), 2023.

\bibitem{zhu2023genimage}
M.~Zhu, H.~Chen, Q.~Yan \emph{et al.}
GenImage: a million-scale benchmark for detecting AI-generated images.
\emph{Advances in Neural Information Processing Systems} (NeurIPS),
Datasets and Benchmarks, 2023.

\bibitem{chen2024drct}
B.~Chen, J.~Zeng, J.~Yang y R.~Yang.
DRCT: diffusion reconstruction contrastive training for synthetic image
detection.
\emph{International Conference on Machine Learning} (ICML), 2024.

\end{thebibliography}

\end{document}
